\documentclass[11pt,letterpaper]{article}

% ===================== Codificación y fuentes =====================
\usepackage{fontspec}
\IfFontExistsTF{Arial}{\setmainfont{Arial}}{\setmainfont{TeX Gyre Heros}}

% ===================== Página y utilidades ========================
\usepackage[left=2.5cm,right=2.5cm,top=2.5cm,bottom=2.5cm]{geometry}
\usepackage{array}
\usepackage{enumitem}
\usepackage{graphicx}
\usepackage{xcolor}
\usepackage{titlesec}
\usepackage{fancyhdr}
\usepackage{hyperref}
\usepackage{booktabs}

% ===================== Etiquetas en español (sin babel) ===========
\renewcommand{\contentsname}{Índice}
\renewcommand{\figurename}{Figura}
\renewcommand{\tablename}{Tabla}
\renewcommand{\refname}{Bibliografía}

% ===================== Escala de grises (sin color) ================
\definecolor{negro}{RGB}{0,0,0}
\definecolor{gris}{RGB}{90,90,90}
\definecolor{grisclaro}{RGB}{150,150,150}
\definecolor{fillgray}{RGB}{175,175,175}

% ===================== Hipervínculos (negro) =======================
\hypersetup{colorlinks=true, linkcolor=negro, urlcolor=negro, citecolor=negro}

% ===================== Estilo de títulos ===========================
\titleformat{\section}{\large\bfseries\color{negro}}{\thesection.}{0.5em}{}
\titleformat{\subsection}{\normalsize\bfseries\color{negro}}{\thesubsection}{0.5em}{}
\titlespacing*{\section}{0pt}{0.75em}{0.25em}
\titlespacing*{\subsection}{0pt}{0.8em}{0.2em}
\makeatletter
\renewcommand*{\l@section}{\@dottedtocline{1}{0em}{1.5em}}
\makeatother

% ===================== Pie de página ===============================
\pagestyle{fancy}
\fancyhf{}
\fancyfoot[L]{\footnotesize\color{gris}Magíster en Ciencia de Datos e Inteligencia Artificial}
\fancyfoot[R]{\footnotesize\color{gris}\thepage}
\renewcommand{\headrulewidth}{0pt}
\renewcommand{\footrulewidth}{0.4pt}

% ===================== Comandos de guía ============================
% \guia ya NO se imprime: el límite de 3 pags (portada+indice+todo+biblio,
% ver "Aspectos formales" del enunciado) no da espacio para texto de ayuda
% en el PDF final. La guía de cada sección queda como comentario "%GUIA:"
% justo antes del \section correspondiente, visible solo en el .tex.
% Antes de entregar: borrar también los \resp que queden.
\newcommand{\resp}{{\footnotesize\itshape\color{fillgray}[\,Desarrollen aquí\,]}\par\vspace{0.3em}}

% Caja de información de la evaluación (escala de grises):

% ===================== Datos del grupo ==============================
\newcommand{\proyTitulo}{Asociación entre características sociodemográficas y supervivencia en el Titanic: un análisis exploratorio}
\newcommand{\proyIntegrantes}{Felipe Díaz Toro, Pablo Ríos Passteni, Guillermo Subiabre Chacón, Ninoska Yévenes Hernández}
\newcommand{\proyDataset}{\texttt{titanic} (repositorio \texttt{seaborn-data}, Waskom, s.\,f.)}
\newcommand{\proyDocente}{Sergio Paraíso Medina}
\newcommand{\proyFecha}{29 de noviembre de 2026}
% \proyRepo: no lo exige el enunciado ni la rúbrica de Sumativa 1 de MCDI503
% (se tomó de la plantilla de MCDI501, que sí lo pedía). Se deja definido por
% si se quiere agregar de vuelta, pero no se imprime en la portada.
\newcommand{\proyRepo}{\url{https://github.com/fdiaztoro/proyecto-grupo4-mcdi503}}

\setlength{\emergencystretch}{2em}
\setlength{\parindent}{0pt}
\setlength{\parskip}{0.25em}

\begin{document}
\pagenumbering{arabic}

\thispagestyle{empty}
\begin{center}
\vspace*{0.3cm}
\includegraphics[height=3.0cm]{logo_unab.png}\\[0.6cm]

{\footnotesize\color{gris}\MakeUppercase{Magíster en Ciencia de Datos e Inteligencia Artificial}}\\[0.15cm]
{\footnotesize\color{grisclaro}MCDI503: Exploración Inteligente para la Ciencia de Datos}\\[1.0cm]

{\Large\bfseries\color{negro} Evaluación Sumativa 1}\\[0.3cm]
{\large\color{negro} Avance Proyecto Fase 1 --- Implementación inicial de EDA reproducible}\\[0.5cm]

{\small\itshape\color{gris} Informe técnico grupal \textperiodcentered\ 20\,\%}
\end{center}

\vspace{0.8cm}

\begin{flushleft}
\setlength{\extrarowheight}{4pt}
\begin{tabular}{@{}>{\bfseries\color{negro}}l p{10.5cm}@{}}
Título del proyecto: & \proyTitulo \\
Integrantes:         & \proyIntegrantes \\
Dataset:             & \proyDataset \\
Docente:             & \proyDocente \\
Fecha:               & \proyFecha \\
\end{tabular}
\end{flushleft}

\vspace{0.6cm}
\vspace{0.3em}
{\small\tableofcontents}

\refstepcounter{section}
\addcontentsline{toc}{section}{\protect\numberline{1}Portada e índice}
\clearpage

\section{Introducción y contextualización}
Los registros del Titanic (1912) permiten explorar la relación entre características de los pasajeros y supervivencia (notebook, sección II). El dataset \texttt{titanic} (Waskom, s.\,f.) reúne 891 pasajeros con variables sociodemográficas (sexo, edad), de viaje (clase, tarifa, puerto) y de acompañantes, con valores faltantes y tipos de variable mixtos, adecuado para un primer ciclo de EDA reproducible. El diagnóstico orienta la limpieza en las siguientes fases.

\section{Objetivo exploratorio, preguntas iniciales y alcance del avance}
\textbf{Objetivo.} Examinar, mediante un EDA descriptivo y reproducible, la asociación entre la supervivencia y las características sociodemográficas y de viaje de los pasajeros (sexo, edad, clase y acompañamiento familiar), revisando antes la calidad de los datos y documentando las decisiones de preparación (notebook, sección II).

\textbf{Preguntas exploratorias iniciales:}
\begin{enumerate}[nosep]
  \item ¿Cómo se distribuyen la edad, la tarifa y la clase del pasaje? ¿Hay valores potencialmente atípicos? (sección IV)
  \item ¿Cambia la supervivencia según el sexo y la clase del pasaje? (sección VI)
  \item ¿Cómo se relaciona la edad con la supervivencia? (sección VI)
  \item ¿Viajar con familiares registrados se asocia con sobrevivir? (secciones IV y VI.3)
  \item ¿Los faltantes de \texttt{age} y \texttt{deck} se concentran en ciertos grupos? (sección V)
\end{enumerate}

\textbf{Alcance de esta fase.} Diagnóstico de calidad, exploración inicial y comparaciones descriptivas. La tarifa se describe en IV y el puerto se usa en el diagnóstico de faltantes. Los supuestos y límites se precisan en la sección 7 de este informe.

% GUIA (12 pts, junto a Sec. 7 -- "Notebook ejecutado y flujo del EDA"): flujo
% general del notebook, referencia a sus secciones, integracion narrativa+codigo,
% correspondencia pregunta <-> seccion donde se responde.
\section{Implementación inicial del EDA en notebook}
El notebook integra celdas narrativas de pregunta, acción y justificación con código, salidas e interpretación. Verifica y revisa los datos (III), describe las distribuciones originales (IV), prepara una tabla de trabajo (V) y contrasta los grupos (VI). El enfoque de EDA mínimo viable del curso (Paraíso, 2026) orienta los resúmenes con denominadores visibles y el registro de decisiones. En IV se usan histogramas, frecuencias y la regla de 1,5 veces el rango intercuartílico (Bruce et al., 2020). Las visualizaciones usan Matplotlib (Hunter, 2007) y Seaborn (Waskom, 2021).

% GUIA (9 pts): que decisiones se tomaron, por que, que resultado intermedio
% produjeron y como se relacionan con la evolucion del analisis en el notebook.
\section{Documentación trazable de decisiones exploratorias}
Se conservan las filas idénticas porque el dataset carece de identificador individual (III.4 y V.2). Se mantienen las edades y cubiertas faltantes y se crean \texttt{age\_missing} y \texttt{deck\_known}. Se excluyen cinco columnas redundantes tras comprobar sus equivalencias, se tipifican las categóricas y se deriva el tamaño familiar. El resultado es \texttt{titanic\_eda} (891 filas, 13 columnas), que usa VI. La exploración de IV usa los datos originales.

% GUIA (9 pts): coherencia datos de entrada -> transformaciones -> resultados;
% organizacion/ejecucion del notebook; flujo sin pasos omitidos ni ocultos.
\section{Reproducibilidad y coherencia técnica}
En III.1 se verifica el SHA-256 del CSV local o descargado antes de leerlo con \texttt{pandas.read\_csv} (The pandas development team, s.\,f.). Los datos originales se conservan y las transformaciones de la copia se validan en V.3. El notebook registra las versiones de ejecución y declara sus dependencias en \texttt{requirements.txt}. Las 17 celdas se ejecutaron desde cero por ambas rutas con el mismo hash y dimensiones.

% GUIA (6 pts): pasos seguidos, supuestos, limitaciones, pendientes para fases
% posteriores.
\clearpage
\section{Registro del proceso, supuestos y limitaciones}
Se verificó y diagnosticó la fuente, se exploraron las distribuciones, se preparó una copia y se compararon los grupos (III--VI). Cada fila se usa como unidad de análisis. Sin un identificador, no puede determinarse si las filas idénticas representan al mismo pasajero. Los contrastes de edad usan 714 de 891 registros y la ausencia de edad varía entre grupos. La cubierta tiene 77,2\,\% de faltantes, concentrados en segunda y tercera clase. Los resultados se circunscriben al dataset y a comparaciones descriptivas. En la siguiente fase se evaluará la sensibilidad al tratamiento de faltantes y se podrán estratificar las comparaciones.

% GUIA (6 pts): principales hallazgos, interpretacion breve, referencias
% explicitas a secciones/tablas/figuras del notebook.
\section{Hallazgos exploratorios iniciales y referencias al notebook}
\begin{enumerate}[nosep,leftmargin=*]
\item Edad mediana de 28 años (714 registros válidos), tercera clase con 491 de 891 pasajeros y 116 tarifas sobre el límite superior de 1,5 veces el rango intercuartílico (IV, tabla descriptiva y gráficos).
\item La tasa observada es mayor entre mujeres en las tres clases. En primera clase, 91 de 94 mujeres (96,8\,\%) y 45 de 122 hombres (36,9\,\%) sobrevivieron (VI.1, tabla por sexo y clase).
\item Las medianas de edad coinciden en 28 años entre quienes sobrevivieron (290) y quienes no (424), con edad conocida (VI.2, tabla y diagrama de caja).
\item La tasa global más alta se observa en grupos familiares de 2 a 4 personas (57,9\,\%, 292 pasajeros). La comparación entre viajar solo o acompañado cambia al separar por sexo (VI.3).
\item En tercera clase se observa un 27,7\,\% de edad faltante y un 97,6\,\% de cubierta faltante (V.1, tabla por clase).
\end{enumerate}

% GUIA (incluido en "Aspectos formales institucionales", 6 pts): minimo 3
% fuentes APA 7a ed. (docente + tecnica + complementaria).
\section{Bibliografía}
Bruce, P., Bruce, A., \& Gedeck, P. (2020). \textit{Practical statistics for data scientists} (2.\textsuperscript{a} ed.). O'Reilly Media.

Hunter, J. D. (2007). Matplotlib: A 2D graphics environment. \textit{Computing in Science \& Engineering}, \textit{9}(3), 90--95. \url{https://doi.org/10.1109/MCSE.2007.55}

Paraíso, S. (2026). \textit{Resumen aplicado del libro: EDA mínimo viable} [Apunte]. Universidad Andrés Bello.

The pandas development team. (s.\,f.). \textit{pandas.read\_csv}. pandas documentation. \url{https://pandas.pydata.org/docs/reference/api/pandas.read_csv.html}

Waskom, M. (s.\,f.). \textit{seaborn-data} [Conjunto de datos]. GitHub. \url{https://github.com/mwaskom/seaborn-data}

Waskom, M. L. (2021). seaborn: Statistical data visualization. \textit{Journal of Open Source Software}, \textit{6}(60), 3021. \url{https://doi.org/10.21105/joss.03021}

\end{document}
